\documentclass[../../main.tex]{subfiles}

\begin{document}
\begin{CJK*}{UTF8}{gbsn}

\section*{Exercise 2}
Consider split conformal prediction 
with data set $(X_1, Y_1), \cdots, (X_{n+2}, Y_{n+2}) 
\in \mathcal{X} \times \mathcal{Y}$ 
with given score 
$s : \mathcal{X} \times \mathcal{Y} \to \mathbb{R}$.
We write $S_i = s(X_i, Y_i)$ for all $i$ for simplicity.
For $j \in \{1,2\}$, the conformal $p$-value is the random variable defined by 

\begin{equation*}
    p_j = \frac{1+\sum_{i=1}^{n}\chi_{ \curly{S_{n+j}\leqslant S_i}}}{n+1}.
\end{equation*}
Given $\alpha \in (0,1)$,
let $\mathcal{C}$ be the conformal interval 
calibrated only on the first $n$ observations.
In other words, for $j \in \{1,2\}$, 
the event $\{Y_{n+j}\in\mathcal{C}(X_{n+j}) \}$
equals to $\{p_j < \alpha\}$.
Assuming that $((X_1, Y_1), \cdots, (X_{n+2}, Y_{n+2}))$ 
is exchangeable, $s$ is a symmetric score function, 
and that the random vector $(p_1,p_2)$
satisfies the 
positive regression dependence on subsets (PRDS) assumption 
given in definition 10.3 of the textbook,
prove that 

\begin{equation*}
    P\{Y_{n+1}\in\mathcal{C}(X_{n+1}),\,
       Y_{n+2}\in\mathcal{C}(X_{n+2})\}\geqslant(1-\alpha)^2
\end{equation*}
without explicitly computing the probability on the left 
hand side of the inequality.

\paragraph{Solution}
First of all, Theorem 3.2 in the textbook already guarantees that,
for $j \in \{1,2\}$:

\begin{equation*}
    P\{Y_{n+j}\in\mathcal{C}(X_{n+j})\} \geqslant 1-\alpha.
\end{equation*}
We focus on showing that 

\begin{equation*}
    P\curly{Y_{n+2}\in\mathcal{C}(X_{n+2}) \mid Y_{n+1}\in\mathcal{C}(X_{n+1}) } 
    \geqslant 1-\alpha.
\end{equation*}
In order to make use of the PRDS condition, 
define a function $g$ on the support of the discrete 
random variable $p_1$ by setting,
for all $t$, that 

\begin{equation*}
    g(t) = P \bracket{ \curly{p_2 > \alpha} \mid p_1 = t}.
\end{equation*}
Note the following facts:

\begin{enumerate}
    \item $g$ is monotonically increasing because 
$g(t) = P \bracket{\curly{(p_1,p_2)\in A} \mid p_1 = t}$
with the set $A$ defined by all $(u,v) \in \mathbb{R}^2$ 
with $v < \alpha$, who is an increasing set.
    \item The average value of $g(p_1)$ on the set 
    $\{p_1 > \alpha\}$ is the desired conditional probability.
    In symbols:

    \begin{equation*}
        \frac{1}{P\{p_1 > \alpha\}} \mathbb{E} \squared{g(p_1); \curly{p_1 > \alpha}}
        = P\curly{Y_{n+2}\in\mathcal{C}(X_{n+2}) \mid Y_{n+1}\in\mathcal{C}(X_{n+1}) }.
    \end{equation*}
    \item The expectation of $g(p_1)$ is the marginal covergae of $(X_{n+2}, Y_{n+2})$.
    In symbols:
    \begin{equation*}
        \mathbb{E} \squared{g(p_1)}
        = P\curly{Y_{n+2}\in\mathcal{C}(X_{n+2})}.
    \end{equation*}
    \item The average value of $g(p_1)$ on the set $\{p_1 > \alpha\}$
    is greater or equal to the average value of $g(p_1)$ on the set 
    $\{p_1 \leqslant \alpha\}$ because $g$ is monotonically increasing.
    That is:

    \begin{equation*}
        \frac{1}{P\{p_1 > \alpha\}} \mathbb{E} \squared{g(p_1); \curly{p_1 > \alpha}}
        \geqslant \frac{1}{P\{p_1 \leqslant \alpha\}} 
        \mathbb{E} \squared{g(p_1); \curly{p_1 \leqslant \alpha}}
    \end{equation*}
\end{enumerate}
The proof is complete once we combine the last three facts 
with the marginal coverage on $(X_{n+2},Y_{n+2})$:

\begin{equation*}
    \begin{split}
    & \tab[0.4cm] 1- \alpha \geqslant P\curly{Y_{n+2}\in\mathcal{C}(X_{n+2})}
    = \mathbb{E} \squared{g(p_1)} 
    \\ & = \frac{P\{p_1 > \alpha\}}{P\{p_1 > \alpha\}} \mathbb{E} \squared{g(p_1); \curly{p_1 > \alpha}}
    + \frac{P\{p_1 \leqslant \alpha\}}{P\{p_1 \leqslant \alpha\}} 
        \mathbb{E} \squared{g(p_1); \curly{p_1 \leqslant \alpha}}
    \\ & \geqslant 
    (P\{p_1 > \alpha\} + P\{p_1 \leqslant \alpha\})
    \frac{1}{P\{p_1 > \alpha\}} \mathbb{E} \squared{g(p_1); \curly{p_1 > \alpha}}
    \\ & = \frac{1}{P\{p_1 > \alpha\}} \mathbb{E} \squared{g(p_1); \curly{p_1 > \alpha}}
    = P\curly{Y_{n+2}\in\mathcal{C}(X_{n+2}) \mid Y_{n+1}\in\mathcal{C}(X_{n+1}) }.
    \end{split}
\end{equation*}
////

\end{CJK*}
\end{document}
